\chapter{Wind turbines}

% --- Content starts here ---

\section{Summary: Ventis -- Wind Farm of Gouy-Lez-Pi\'etons}

\subsection{Company Overview}

Ventis S.A. is a Belgian wind energy developer founded in the early 2000s by brothers Benoît
and Pierre MAT, headquartered in Tournai. The company manages wind energy projects entirely
in-house -- from development and permitting through construction and long-term operation --
operating with a local and sustainable philosophy in Wallonia and France. Its internal team
covers all key disciplines: development, construction, operation, administration, and
communication.

Between 2022 and 2025, Ventis installed 25\% of all new wind capacity commissioned in Wallonia,
recently becoming the first independent wind energy producer in the region. The company currently
operates approximately 90 wind turbines across Belgium and France, representing a total installed
capacity of around 250\,MW, of which nearly 200\,MW corresponds to its own ownership share.

In Belgium, Ventis operates wind farms across the provinces of Hainaut, Brabant wallon, and
Li\`ege, with sites in operation such as Dour--Qui\'evrain--Hensies (9 turbines),
Tournai-Antoing-Brunehaut (10 turbines), La Louvi\`ere (5 turbines), and Qu\'evy-Mons
(7 turbines), among others. Several projects are also authorised or under development,
including repowering operations and new greenfield sites. In France, Ventis operates farms
in Longueville-sur-Aube (Aube, 11 turbines) and Gorges-Gonfreville (Manche, 9 turbines),
with additional projects under development in Pas-de-Calais and Aisne.

\subsection{The Gouy-Lez-Pi\'etons Project}

\subsubsection{Permitting History}

The Gouy-Lez-Pi\'etons wind farm illustrates the complexity of the Belgian permitting process
for wind energy. Ventis submitted its initial permit application in March 2014. The first
decision, issued in August 2014 by the delegated and technical officers, was a refusal based
on an unfavourable opinion from Belgocontrol (the Belgian air navigation service provider).
Ventis appealed, obtained a conditional opinion from Belgocontrol, and was granted the permit
on appeal by Minister Carlo Di Antonio in January 2015.

However, the municipality of Pont-\`a-Celles lodged a further appeal in March 2015, citing
high landscape impact. The Council of State annulled the permit in March 2017 on the grounds
that alternative locations had not been sufficiently examined. Ventis resubmitted the permit
application in May 2019, providing a supplemental impact study. The permit was re-granted by
the Minister in July 2019, followed by yet another appeal from the municipality in September
2019, this time contesting the examination of alternatives. Finally, in November 2022, the
Council of State cleared the permit of all appeals, confirming that no alternative locations
were available. From first application to final clearance, the permitting process spanned
nearly nine years.

\subsubsection{Technical Characteristics}

The wind farm comprises 8 turbines arranged across agricultural land near Gouy-Lez-Pi\'etons,
connected by underground MV cables and internal access roads. The main technical characteristics
are as follows:

\begin{itemize}
  \item \textbf{Wind turbine model:} Nordex N131
  \item \textbf{Nominal power per turbine:} 3.6\,MW
  \item \textbf{Total installed capacity:} 28.8\,MW
  \item \textbf{Total height:} 150\,m
  \item \textbf{Rotor diameter:} 131\,m
  \item \textbf{Grid connection:} DSO ORES, 10.8\,kV, 25\,MW injection capacity
  \item \textbf{Annual energy production:} $\approx$57\,GWh
  \item \textbf{Capacity factor:} $\approx$26\%
  \item \textbf{Households supplied:} $\approx$14\,400 (yearly average consumption)
\end{itemize}

\subsubsection{Ownership Structure}

Six of the eight turbines are owned by the project company GOUYNRG, while the remaining two
are owned by ENGIE, which was a competing developer during the development phase. The GOUYNRG
shares are distributed as follows: 70\% held by Ventis, 15\% by Greenplug (a development
partner), and 15\% by local public and citizen participation -- reflecting Ventis's commitment
to community involvement and local acceptability.

\subsection{Project Development}

\subsubsection{Wind Turbine Sizing}

The dimensions of the wind turbines are constrained by a combination of biological and
air traffic factors. The minimum ground clearance (distance between the blade tip at its lowest
point and the ground) is determined by biological constraints, particularly the presence of
bat species requiring a safety margin above the tree canopy. The maximum blade tip height is
set by air traffic constraints enforced by Belgocontrol. Together, these two values define
the admissible rotor diameter (their difference) and, consequently, the tower height. The
wind resource assessment further informs the optimal rotor size, since power extracted from
the wind scales with the swept area and the cube of wind speed:
\[
  P_{\text{wind}} = \frac{1}{2} \rho A v^3 \quad [\text{W}]
\]
where $\rho$ is air density (approximately 1.2\,kg/m$^3$ at 15$^\circ$C), $A$ is the rotor
swept area, and $v$ is wind speed.

\subsubsection{Turbine Layout and Wake Effects}

The number and spatial distribution of turbines is governed by the available grid injection
capacity (25\,MW with ORES), land acquisition constraints, local setback distances from
housing, natural environments, and infrastructure, and wake effect losses. Wake effects occur
when a downstream turbine operates in the turbulent, lower-energy airflow created by an
upstream turbine. The prevailing wind rose at the site (predominantly south-southwest to
west-southwest) directly influences the optimal turbine spacing and alignment to minimise
these losses.

A constraint map overlaying setback distances from housing zones, motorways, high-voltage
lines, hertzian beams, forests, and Natura 2000 areas was produced during the development
phase to define the feasible turbine positions.

\subsection{Turbine Technology}

\subsubsection{Turbine Selection Process}

Three candidate turbine models were evaluated: the Vestas V126 (3.45\,MW), the Nordex N131
(3.60\,MW), and the Siemens Gamesa SG-132 (3.47\,MW). Selection criteria included permit
height limits (total height and ground clearance), rotor size, nominal power, noise output,
site access feasibility, wind resource studies, and site assessment reviews.

A key finding during the site assessment was that turbine E3 (and to a lesser extent E6)
faced a 15\% energy curtailment constraint due to the proximity of a small woodland, reducing
the net annual energy production of the SG-132 option to only 49\,130\,MWh/year -- lower than
the N131's 57\,688\,MWh/year. Combined with a competitive price-to-performance ratio and no
turbulence-related power limitation, the Nordex N131 was selected.

\subsubsection{Generator Architectures}

Two main electrical architectures are used in the wind industry today:

\paragraph{Doubly-Fed Induction Generator (DFIG).}
Used by manufacturers such as Nordex, Vestas, and Siemens Gamesa (onshore). The rotor is
connected to the grid through a partial-scale power converter (handling only 25--35\% of
rated power), while the stator is connected directly. This results in a smaller, cheaper
converter and high efficiency at rated power. However, the direct stator--grid coupling makes
the machine more sensitive to grid voltage fluctuations, and the gearbox and slip rings
require more frequent maintenance.

\paragraph{Gearless Direct Drive (Permanent Magnet Synchronous Generator -- PMSG).}
Used by Enercon and Siemens Gamesa (offshore). The rotor shaft drives the generator directly,
eliminating the gearbox. A full-scale power converter decouples the generator entirely from
the grid, enabling superior grid performance (fault ride-through, precise active and reactive
power control). The absence of a gearbox reduces maintenance requirements and potentially
extends the turbine's operational lifetime. The trade-off is a larger, heavier, and more
expensive full-power converter chain, and more complex fault diagnostics. Enercon's design
has evolved over generations: older models (e.g.\ E82) placed the rectifier in the nacelle
and the inverter at the tower base via a DC link, while current models (E115) integrate the
full conversion chain in the nacelle, and newer large models (E175) consolidate all power
electronics in a compact nacelle structure.

\subsection{Construction}

\subsubsection{Logistics and Transport}

Wind turbine transportation is one of the most challenging aspects of onshore wind projects
in densely populated regions such as Belgium. Components -- blades (approximately 65\,m),
tower sections, and nacelles (approximately 13\,m long, 4\,m wide) -- are manufactured in
Europe and transported by sea to the port of Antwerp. A local specialised transport company
then takes over responsibility for the convoy from the port to the site, including liability
for damage or blockages.

The transport company conducts detailed route surveys, produces turn simulations using
dedicated software, and performs dummy runs (typically at night) to validate the route before
actual component delivery. Road adaptations may include removing road signs, applying steel
plates on roundabout islands, and using high-density polyethylene (PEHD) ground protection
plates to extend turning radii on agricultural land. From the road transition point (RTP)
onwards, Ventis assumes full responsibility for road design, land agreements, and any delays
or damage.

\subsubsection{Civil Works (Balance of Plant)}

Access roads and crane platforms must meet strict bearing capacity requirements: up to
12\,t per axle and a total crane load of up to 137\,t. Road layers are designed using
finite element modelling (FEM) based on geotechnical reports and OEM specifications, and
validated on-site by plate load tests (EV1/EV2 compaction ratios). Water management
(drainage, avoiding mud and icing) is also a critical design criterion.

\paragraph{Foundations.}
Two main categories of foundation design are used: slab (or gravity) foundations and pile
foundations. The choice depends on geological conditions (pile foundations are unsuitable
for karstic limestone), available land area, groundwater level (foundations are elevated to
reduce buoyancy and increase hub height), and cost (depth to competent rock is a key driver).
A geotechnical campaign -- comprising penetration tests, core drilling, electrical
resistivity tomography, and piezometer installation -- is mandatory to characterise the
subsoil. The Gouy-Lez-Pi\'etons foundations have a diameter of 19.2\,m, compared to
17.6\,m for a comparable 2\,MW turbine from 2007, reflecting the increased loads of modern
larger machines.

Gravity foundation variants include direct soil replacement and soil improvement techniques
(ballasted columns, rigid inclusions, or mixed inclusions). Pile foundations may use driven
piles (steel or concrete, impact-driven) or bored piles (various auger techniques).
Concreting of the slab is performed in a single continuous pour, and the foundation is
then backfilled and compacted.

\subsubsection{Electrical Infrastructure}

MV cables (24\,kV, aluminium, cross-sections of 150 or 240\,mm$^2$) are buried at a
depth of approximately 1.3\,m, laid in a bed of sand with PE protection slabs and warning
ribbons, alongside HDPE conduits for fibre-optic communication cables.

The on-site electrical cabin contains two distinct sections. The high-voltage (HV) side
houses ten SF6 switchgear cells: three for ENGIE (departure, metering, generator protection),
four for ORES (two arrival feeders, reference voltage, ORES protection), and three for
Ventis/GOUYNRG (generator protection, metering, departure), plus an auxiliary transformer.
The low-voltage (LV) side contains the SCADA controllers for each owner (Nordex CWE units),
an aggregator box for curtailment management based on electricity prices, a NORTEC unit for
automated shadow flicker and bat curtailment, a Phoenix UPS, internet routers with fixed IPs
and firewalls, and a 4G backup system.

\subsubsection{Turbine Assembly and Commissioning}

Assembly involves the positioning of a main crawler crane (LTR1100) and an assist crane
(LTM1650-8.1) according to OEM lifting plans. Tower sections, nacelle, hub, and blades are
lifted and assembled sequentially. Commissioning includes electrical completion, fibre-optic
cabling, functional test runs, grid energisation, and formal inspections by the DSO and a
certified independent body (e.g.\ TÜV, Vinçotte, Bureau Veritas). A punch list process
formalises the handover between the OEM and the owner.

\subsection{Operation and Maintenance}

Operation of the wind farm encompasses preventive maintenance, corrective maintenance,
remote monitoring, mandatory regulatory inspections, and environmental monitoring.

Preventive maintenance includes periodic bolt re-torquing (tower flanges, nacelle, blade
roots), regreasing of gearboxes, pitch bearings, and yaw rings, visual inspections of blades,
nacelle components and tower internals, functional testing of safety systems, and planned
replacement of wear parts. A condition monitoring system (CMS) using vibration sensors on
the main bearing and gearbox allows early detection of anomalies such as under-greased
bearings, enabling intervention before catastrophic failure.

Corrective maintenance covers unplanned repairs such as blade damage from lightning strikes
(repaired by rope-access technicians) and major component replacements such as generator
swaps, which require a dedicated service crane.

Remote monitoring is performed via the Farsight SCADA platform, which provides real-time
dashboards, power curve analysis, wind analysis, alarm management, and reporting for each
turbine across the portfolio.

Environmental monitoring obligations include post-construction bat and bird mortality surveys
(using trained dogs and regular transect walks), shadow flicker management (automated
curtailment via NORTEC when a receptor exceeds the permitted annual shadow hours), noise
measurements at neighbouring dwellings, and ice detection on blades.

\subsection{Grid Constraints and Power Offtake Optimisation}

\subsubsection{Curtailments}

The 25\,MW injection capacity agreed with ORES limits the farm's ability to export its full
28.8\,MW installed capacity at all times. Additionally, curtailments are triggered when day-ahead
or intraday electricity market prices turn negative, since injecting power at negative prices
generates a direct financial loss. Environmental curtailments (bats, shadow, noise) add further
stops, typically representing 4\% of gross annual production at Gouy.

\subsubsection{Grid Congestion Context}

The Belgian distribution grid is under significant stress: renewable energy connection requests
to ORES increased by more than a factor of four between recent years, and the total demand
connected or pending connection is growing rapidly (from approximately 350\,MW measured load
to projections above 500\,MW by 2025), driven by battery storage, EV fast charging, and
industrial electrification. The grid is aging and was designed for a unidirectional power flow
from large central plants to consumers -- the opposite of distributed renewable generation.

\subsection{Future Developments}

Ventis identifies four main optimisation vectors to improve the economic viability of its
portfolio under increasing grid congestion and market volatility:

\paragraph{Behind-the-meter direct connection (private PPA).}
A direct electrical line between a wind farm and a nearby industrial customer bypasses the
distribution grid entirely, eliminating network tariffs. Both parties benefit: the producer
receives a higher price than the spot market, and the industrial customer pays less than the
regulated tariff. This arrangement also tends to facilitate permitting, as industrials
committing to green energy consumption are viewed favourably. Battery storage can further
improve the self-consumption rate.

\paragraph{Repowering.}
End-of-life turbines are replaced by fewer but larger and more powerful machines, making use
of the existing grid connection, partially secured land, and an already-integrated landscape.
Modern turbines deliver significantly more energy per unit cost than those installed 20 years
ago. Permitting is generally more straightforward, as the visual impact on the landscape is
already accepted by the local community and the turbine count decreases.

\paragraph{Battery Energy Storage Systems (BESS).}
BESS co-located with wind farms can absorb energy that would otherwise be curtailed (due to
grid limits or negative prices) and re-inject it later. They also enable participation in
ancillary service markets (frequency containment reserve FCR, automatic frequency restoration
reserve aFRR, imbalance settlement) and support peak-shaving strategies. One concrete project
under development involves a 2\,MW / 4\,MWh BESS at a 10.8\,MW wind farm (PEDG) with an
8.8\,MW injection limit.

\paragraph{EV fast-charging stations.}
Wind turbines can supply EV chargers through a direct private line, providing traceable
renewable electricity for mobility without loading the public grid. This creates a win-win
situation for the producer (better selling price) and the consumer (cheaper green electricity),
reduces the need for municipalities to fund public charging infrastructure, and increases
local acceptance of associated storage or wind projects. A concept project (MWPk) envisions
approximately 10 fast chargers near a motorway junction with a potential of 50\,000 vehicles
per day, powered by a 9.4\,MW wind farm with no impact on the local distribution grid.

\paragraph{Wind/solar hybridisation.}
Solar PV generation is strongly complementary to wind in Belgium: wind output peaks in winter
and at night, while solar generation peaks in summer and midday. Combining both sources at a
shared grid connection point improves the utilisation of the grid capacity and the associated
BESS, reduces the net cost of grid connection for the solar project, and creates opportunities
for agrivoltaic development beneficial to local farmers. A project under study (EAE) combines
a 25\,MW wind farm with a 700\,kWp solar installation and a 3.5\,MW / 7\,MWh BESS.